%\documentclass[12pt,letterpaper]{article}

%%%
%% Required packages:
%%
\usepackage{etex}
\usepackage{amsmath}
\usepackage{amssymb}
\usepackage{amstext}
\usepackage{amsthm}
\usepackage{fancyhdr,fancybox}
\usepackage{mathrsfs} % need for \mathscr command
\usepackage{multicol}
\usepackage[normalem]{ulem}
%\usepackage{pictex,pictexplus}
% \usepackage{pictexwd}
\usepackage{rotating}
\usepackage{url}
\usepackage{xfrac}
\usepackage{booktabs}
\usepackage{color,soul}
\usepackage{comment}

\usepackage{hyperref}
\hypersetup{
    colorlinks=true,     % false: boxed links; true: colored links
    linkcolor=blue,      % color of internal links
    citecolor=blue,      % color of links to bibliography
    filecolor=blue,      % color of file links
    urlcolor=blue        % color of external links
}


\usepackage{tikz}
\usetikzlibrary{backgrounds,calc}

%%
%% Common formatting commands
%%
\setlength{\parindent}{0in}
\setlength{\textwidth}{7in}
\setlength{\evensidemargin}{-0.25in}
\setlength{\oddsidemargin}{-0.25in}
\setlength{\parskip}{.5\baselineskip}
\setlength{\topmargin}{-0.5in}
\setlength{\textheight}{9in}

%               Problem and Part
\newcounter{problemnumber}
\newcounter{partnumber}
\newcounter{subpartnumber}
\newcommand{\Problem}{%
  \refstepcounter{problemnumber}%
  \setcounter{partnumber}{0}%
  \setcounter{subpartnumber}{0}%
  \item[\makebox{\hfil\textbf{\theproblemnumber.}\hfil}]%
  }
\newcommand{\InProblem}[2][1.6]{%
  \refstepcounter{problemnumber}%
  \setcounter{partnumber}{0}%
  \setcounter{subpartnumber}{0}%
  \textbf{\theproblemnumber.}\ \ \parbox[t]{#1in}{#2}%
}
\newcommand{\InSmallProblem}[1]{\InProblem[1.05]{#1}}
\newcommand{\NoProblem}[1]{%
  \mbox{\hfil\textbf{#1}\hfil}%
  }
\newcommand{\UnProblem}{%
  \refstepcounter{problemnumber}%
  \setcounter{partnumber}{0}%
  \setcounter{subpartnumber}{0}%
  \mbox{\hfil\textbf{\theproblemnumber.}\hfil}%
  }
\newcommand{\InFlexProblem}[1]{%
  \refstepcounter{problemnumber}%
  \setcounter{partnumber}{0}%
  \setcounter{subpartnumber}{0}%
  \textbf{\theproblemnumber.}\ \ #1%
}
%%  Part:
\newcommand{\Part}{%
  \renewcommand{\thepartnumber}{(\alph{partnumber})}%
  \refstepcounter{partnumber}
  \setcounter{subpartnumber}{0}%
  \item[(\alph{partnumber})]%
}
\newcommand{\InPart}[2][1.75]{%
  \renewcommand{\thepartnumber}{(\alph{partnumber})}%
  \refstepcounter{partnumber}%
  \setcounter{subpartnumber}{0}%
  (\alph{partnumber})\ \ \parbox[t]{#1in}{#2}%
}
\newcommand{\UnPart}{%
  \refstepcounter{partnumber}%
  \renewcommand{\thepartnumber}{(\alph{partnumber})}%
  \setcounter{subpartnumber}{0}%
  (\alph{partnumber})%
}
\newcommand{\InLargePart}[1]{\InPart[3.05]{#1}}
\newcommand{\InFullPart}[1]{\InPart[6.2]{#1}}
\newcommand{\InSmallPart}[1]{\InPart[1.05]{#1}}
%% SubPart:
\newcommand{\SubPart}{%
  \refstepcounter{subpartnumber}%
  \item[(\roman{subpartnumber})]%
  }
\newcommand{\InSubPart}[2][2.25]{%
  \stepcounter{subpartnumber}%
  (\roman{subpartnumber})\ \ \parbox[t]{#1in}{#2}%
}
\newcommand{\InSmallSubPart}[1]{\InSubPart[1.5]{#1}}
\newcommand{\InTinySubPart}[1]{%
  \stepcounter{subpartnumber}%
  (\roman{subpartnumber})\ \ \makebox{#1}%
}
\newcommand{\InMySubPart}[2][2.25]{%
  \stepcounter{subpartnumber}%
  \parbox[t]{#1in}{\makebox[0.3in][l]{(\roman{subpartnumber})}\ \ {#2}}%
}
\newcommand{\TrueFalse}{\textbf{T}\ \ \textbf{F}\ \ }
\newcommand{\TFPart}[1]{% 
  \stepcounter{partnumber}%
  \setcounter{subpartnumber}{0}%
  \item[(\alph{partnumber})] \TrueFalse\parbox[t]{5.5in}{#1}}

\newcommand{\Points}[1]{(#1 points) \ }
\newcommand{\Point}[1]{(#1 point) \ }


%% For Answers:
\newcounter{answernumber}
\newcommand{\Answer}{%
  \refstepcounter{answernumber}%
  \setcounter{partnumber}{0}%
  \setcounter{subpartnumber}{0}%
  \item[\makebox{\hfil\textbf{\theanswernumber.}\hfil}]%
  }


% Setting up the headers for the first page
%\chead{} % will default to what is typed in the file.
\fancyhead[L]{\textsc{Math 22a}}
\fancyhead[R]{\textsc{Fall 2024}}
\fancyfoot[R]{
	\vspace*{\fill}\footnotesize\textcopyright{} \emph{Harvard Math 22a}	
}
\renewcommand{\headrulewidth}{0pt}

%\fancyhead[C]{}
%	\chead{}	
%	\lhead{}
%	\rhead{}
%	\cfoot{}


% Putting copyright on the footer of each page
%\fancypagestyle{secondpage}{
%	\fancyhead[C]{TESTing again} % change this to be blank
%		%\chead{}	
%	\fancyhead[L]{bears} % change this to be blank
%	\fancyhead[R]{dogs} % change this to be blank
%	\fancyfoot[R]{ %keeps the footer the same
%		\vspace*{\fill}\footnotesize\textcopyright{} \emph{Harvard Math 22a}	
%	}
%}
%\renewcommand{\headrulewidth}{0pt}
%\pagestyle{fancy}
\pagestyle{empty}


%\lhead{}
%\rhead{}
%\rfoot{\vspace*{\fill}\footnotesize\textcopyright{} \emph{Harvard Math 22a}}
%\renewcommand{\headrulewidth}{0pt}
%\pagestyle{fancy}




%%
%% Common shortcut commands:
%%
\newcommand{\ds}{\displaystyle}
\newcommand{\sgn}{\operatorname{sgn}}
\renewcommand{\Pr}{\operatorname{P}}
\newcommand{\CPr}[3][{}]{\Pr_{\text{#1}}\left(#2\,|\,#3\right)}

\newcommand{\C}{\mathbb{C}}
\newcommand{\R}{\mathbb{R}}
\newcommand{\Z}{\mathbb{Z}}
\newcommand{\N}{\mathbb{N}}
\newcommand{\Q}{\mathbb{Q}}
\newcommand{\D}{\mathbb{D}}

\newcommand{\rref}{\operatorname{rref}}
\newcommand{\im}{\operatorname{im}}
\newcommand{\rank}{\operatorname{rank}}
\newcommand{\diag}{\operatorname{diag}}
\newcommand{\Span}{\operatorname{span}}
%\renewcommand{\vec}[1]{\mathbf{#1}}
\newcommand{\charpoly}{\mathscr{P}}
\newcommand{\nicefrac}[2]{\sfrac{#1}{#2}} % using xfrac package
\newcommand{\topology}{\mathcal{T}}
\newcommand{\Inn}{\operatorname{Inn}}

\newcommand{\blankbox}[2]{\fbox{\rule{#1}{0in}\rule{0in}{#2}}}

\newenvironment{myindentpar}[1]%
 {\begin{list}{}%
         {\setlength{\leftmargin}{#1}
          \setlength{\rightmargin}{\leftmargin}}%
         \item[]%
 }
 {\end{list}}
\newtheorem{theorem}{Theorem}
\newtheorem*{NStheorem}{Theorem}
\newtheorem{cor}[theorem]{Corollary}
\newtheorem*{claim}{Claim}
\newtheorem{defn}{Definition}

\newenvironment{amatrix}[1]{%
  \left(\begin{array}{@{}*{#1}{c}|c@{}}
}{%
  \end{array}\right)
}

%--------grstep
% For denoting a Gauss' reduction step.
% Use as: \grstep{\rho_1+\rho_3} or \grstep[2\rho_5 \\ 3\rho_6]{\rho_1+\rho_3}
\newcommand{\grstep}[2][\relax]{%
   \ensuremath{\mathrel{
       {\mathop{\longrightarrow}\limits^{#2\mathstrut}_{
                                     \begin{subarray}{l} #1 \end{subarray}}}}}}
\newcommand{\swap}{\leftrightarrow}



\chead{\Large\textbf{Proofs \& Counterexamples, \fbox\S 4,  5,  9}}% 

%\begin{document}
\thispagestyle{fancy}

\begin{enumerate}
  \Problem Prove or give a counterexample to each of the following
  ``existential statements''. 
  \begin{enumerate}
    \Part  There exists a prime number $p$ with $90<p<105$. 
    \vfill

    \Part There exists a real number $r$ with $r^2+2r+2\le 0$.
    \vfill

    \Part There exists a sequence of 100 consecutive composite numbers.
    \vfill
  \end{enumerate}

  \Problem Prove or give a counterexample to each of the following
  ``universal statements''.
  \begin{enumerate}
    \Part For every real number $r$, $r^2 + 2r+1\ge0$. 
    \vfill

    \Part For all positive integers $n$, the number $\frac{n(n+1)}{2}$ is an integer.
    \vfill

    \Part Every prime number is odd.
    \vfill
  \end{enumerate}

\end{enumerate}
\newpage

Prove or give a counterexample to each of the following statements.
\begin{enumerate}
  \Problem Let $a$ and $b$ be integers. If $a\,\Big{|}\,b$, then
  $a\,\Big{|}\,3b^3-4b^2+18b+2a$.  
  \vfill

  \Problem Let $m$ and $n$ be integers.  If $m$ and $n$ have different
  parity, then their sum is odd.  
  \vfill

  \Problem Let $x$ be a real number.  If $x^2+3x<0$, then $x<0$.
  \vfill

  \Problem (Careful!) Let $a$ and $b$ be integers.  If $a^2(b^2-2b)$
  is odd, then $a$ and $b$ must both be odd.

  \vfill
\end{enumerate}


%\end{document}
\begin{comment}

\newpage
\chead{\Large\textbf{Proofs \& Counterexamples -- Answers / Solutions}}%
\lhead{}
\rhead{}
\thispagestyle{fancy}

\begin{enumerate}
  \Answer
  \begin{enumerate}
    \Part True. Take $p=97$ (or 101, 103).

    \Part This is false, since $r^2+2r+2 = (r^2+2r+1)+1 = (r+1)^2+1
    \ge 1$.

    \Part This is true, but how can we find it? The simplest thing
    would be to find a number $N$ where $N$, $N+1$, $N+2$, and so on
    are divisible, in turn, by $2$, $3$, $4$, and so on.  (Notice that
    we skipped $1$ -- \emph{every}\ number is divisible by $1$.)
    Maybe instead we should find a number $N$ so that $N+2$ is
    divisible by $2$, $N+3$ is divisible by $3$, and so on up to
    $N+101$ is divisible by $101$ (so $2$ through $101$ give $100$
    numbers). In general we want $N+k$ to be divisible by $k$ for all
    $k$ from $2$ through $101$.  This gives us our answer: just take
    $N$ to be a number that is divisible by every number from $2$
    through $101$.  The easiest number (although certainly not the
    smallest) with this property is $N=101!$.  So the numbers
    \begin{equation*}
      101!+2,\ 
      101!+3,\ 
      101!+4,\ 
      \cdots,\ 
      101!+101
    \end{equation*}
    are all composite.
  \end{enumerate}

  \Answer
  \begin{enumerate}
    \Part This is true, since $r^2+2r+1 = (r+1)^2 \ge 0$.

    \Part This is true, and we'll prove it in a number of ways:
    \begin{itemize}
    \item Directly. If $n$ is even, then $n=2k$ for some positive
      integer $k$ and $\frac{n(n+1)}{2} = k(2k+1)$ is an integer.  If
      $n$ is odd, then $n=2k-1$ for some positive integer $k$ and
      $\frac{n(n+1)}{2} = (2k-1)k$ is an integer.

    \item One fact we'll see in the next few days is that
      \begin{equation*}
        1 + 2 + \cdots + n = \frac{n(n+1)}{2}.
      \end{equation*}
      Written more succinctly, this is $\displaystyle\sum_{k=1}^{n} k
      = \frac{n(n+1)}{2}$.  Thus this number is the sum of the first
      $n$ integers, and so must be an integer itself.

    \item Another fact we'll see in the coming days is that
      $\frac{n(n+1)}{2}$ can also be written as $\binom{n+1}{2}$
      (read ``$n+1$ choose $2$'' and called a \emph{binomial
        coefficient}).  This is perhaps easiest to describe as
      \emph{the number of ways to pick a set of $2$ objects from a
        collection of $n+1$ different objects}.  This ``number of
      ways'' is clearly an integer. 

    \end{itemize}



    \Part Well, no: $p=2$ is prime and even.

  \end{enumerate}

  \Answer We provide a direct proof.  If $a\,\big|\,b$, then $b=ka$
  for some integer $k$.  Then
  \begin{equation*}
    3b^3 - 4b^2 + 18 b  + 2a
    = 3k^3 a^3 - 4k^2a^2 + 18ka + 2a
    = \left( 3k^3a^2 - 4k^2a + 18k + 2 \right) a.
  \end{equation*}
  That is, $a\,\big|\,\left(3b^3-4b^2+18b+2a\right)$.

  \Answer If $m$ is even and $n$ is odd, then we can write $m=2k$ and
  $n=2\ell+1$ for integers $k$ and $\ell$. Thus $m+n = 2(k+\ell)+1$ is
  odd.  
  \medskip

  What if $m$ were odd and $n$ even?  Then the argument would be
  exactly the same, except with the roles of $m$ and $n$ reversed.
  Usually we don't repeat the same argument, but instead note that it
  doesn't really matter (for the argument) which one is even and which
  is odd.  We typically begin our proof with something like ``Without
  loss of generality, assume $m$ is even and $n$ is odd.''  This means
  that we could have assumed the opposite parities, but it doesn't
  really matter and the proof is the same.

  \Answer We instead prove the contrapositive, which is equivalent: If
  $x\ge0$, then $x^2+3x\ge 0$.  This statement is nearly obvious: both
  $x^2$ and $3x$ are non-negative, so their sum is as well: $x^2+3x\ge
  0$. 

  \Answer We'll prove the contrapositive: \emph{if at least one of $a$
    or $b$ is even, then $a^2(b^2-2b)$ is even}.  There are two cases:
  If $a$ is even, then $a = 2k$ and $a^2(b^2-2b) = 2\left[
    2k^2(b^2-2b) \right]$ is even.  If $b$ is even, then $b=2m$ and
  $a^2(b^2-2b) = a^2(4m^2 - 4m) = 2\left[ a^2(2m^2-2m) \right]$ is
  even.  Of course, it could be the case that both $a$ and $b$ are
  even, so we could use either argument.

\end{enumerate}

\end{comment}
%\end{document}


%%% Local Variables: 
%%% mode: latex
%%% tex-master: t
%%% end: 
